\section{Questão 3}

\subsection{Enunciado}

Considere o seguinte conjunto de dados:
\begin{verbatim}
data = {
    'Matemática': [85, 92, 78, 96, 88],
    'Português': [7, 5, 8, 9, 6]
}
\end{verbatim}
É conhecido que as notas de Matemática variam de 0 a 100, enquanto as notas de Português variam de 0 a 10. Aplique a normalização Min-Max e a padronização Z-Score para as notas em ambas as disciplinas, de modo que elas estejam na mesma escala e distribuição.

\subsection{Solução}

As notas de Matemática pertencem à faixa de 0 a 100 e as de Português à faixa de 0 a 10. Os valores brutos, por isso, não compartilham escala. A normalização Min-Max da aula reescala cada coluna para o intervalo \([0, 1]\):
\[
z_{\mathrm{norm}} = \frac{x - x_{\min}}{x_{\max} - x_{\min}}.
\]
\(x_{\min}\) e \(x_{\max}\) são o mínimo e o máximo observados na própria coluna, como no exemplo numérico da aula. A padronização Z-score coloca cada coluna com média zero e desvio padrão unitário:
\[
z = \frac{x - \mu}{\sigma}.
\]
O script ajusta \texttt{MinMaxScaler} e \texttt{StandardScaler}. O desvio é o populacional, o mesmo do exemplo da aula. Em Matemática, o mínimo é 78, o máximo é 96, a média é 87,8 e o desvio é 6,145. Em Português, o mínimo é 5, o máximo é 9, a média é 7 e o desvio é 1,414.

\subsection{Resultado}

Depois das duas transformações, as disciplinas passam a ocupar o intervalo \([0, 1]\) e, na padronização, a mesma distribuição (média 0 e desvio 1). A Tabela~\ref{tab:q3-escalas} reúne os valores. O arquivo \path{scripts/results/q3_escalas.csv} guarda a tabela.

\begin{table}[H]
    \centering
    \caption{Notas originais, normalização Min-Max e padronização Z-score.}
    \label{tab:q3-escalas}
    \footnotesize
    \setlength{\tabcolsep}{4pt}
    \begin{tabular}{@{}rrrrrr@{}}
        \toprule
        \shortstack{Mat.\\orig.} &
        \shortstack{Port.\\orig.} &
        \shortstack{Mat.\\Min-Max} &
        \shortstack{Port.\\Min-Max} &
        \shortstack{Mat.\\Z} &
        \shortstack{Port.\\Z} \\
        \midrule
        85 & 7 & 0,389 & 0,500 & $-0{,}456$ & $0{,}000$ \\
        92 & 5 & 0,778 & 0,000 & $0{,}683$ & $-1{,}414$ \\
        78 & 8 & 0,000 & 0,750 & $-1{,}595$ & $0{,}707$ \\
        96 & 9 & 1,000 & 1,000 & $1{,}334$ & $1{,}414$ \\
        88 & 6 & 0,556 & 0,250 & $0{,}033$ & $-0{,}707$ \\
        \bottomrule
    \end{tabular}
\end{table}
